\section*{Refined Thesis Logic for Fast Area-Weighted Convex Hull Peeling}

\subsection*{Problem framing}
Let $P$ be a set of $n$ points in the plane. The area-weighted peeling process repeatedly removes the hull point whose deletion gives the largest decrease in convex hull area. For a point $p \in P$, define its sensitivity as
\[
\sigma(p) = \mathrm{area}(\mathrm{CH}(P)) - \mathrm{area}(\mathrm{CH}(P \setminus \{p\})).
\]
Only points on the current convex hull can have positive sensitivity, so each peeling step selects
\[
\operatorname*{argmax}_{p \in \mathrm{CH}(P)} \sigma(p).
\]
The process is a greedy heuristic for area-based outlier detection: the first $k$ peeled points are interpreted as the $k$ most extreme outliers.

\subsection*{Relation to exact $k$-removal}
The natural exact optimization problem asks for a set $S$ of $k$ points minimizing the area of $\mathrm{CH}(P \setminus S)$. This exact objective is combinatorial and significantly more expensive than repeated greedy $1$-peels. The weighted peeling algorithm therefore occupies a practical middle ground:
\begin{itemize}
    \item it uses the same geometric objective as the exact problem,
    \item it preserves the interpretation of outliers as points that strongly determine hull area,
    \item and it avoids the high running time of exact $k$-removal algorithms.
\end{itemize}
The tradeoff is that greedy peeling can differ from the optimal joint $k$-removal when nearby outliers mask one another.

\subsection*{Core structural observation}
For a hull point $u$, the change in hull area after deleting $u$ depends only on local hull structure and on the points that would become exposed after removing $u$. These points are the active points of $u$. The crucial structural fact behind the fast algorithm is that although a single peel may trigger many local changes, the total number of times points become newly active over the full process is linear.

\subsection*{Algorithmic decomposition}
The fast algorithm can be understood as maintaining three coupled structures:
\begin{enumerate}
    \item the outer convex layer $L_1$,
    \item the second convex layer $L_2$,
    \item and the remaining interior points in a dynamic hull structure.
\end{enumerate}
At each peel:
\begin{enumerate}
    \item choose the point of maximum sensitivity from $L_1$,
    \item remove it,
    \item restore $L_1$ and $L_2$ using tangent and extreme-point queries,
    \item update the sensitivities of the affected neighboring hull points.
\end{enumerate}
The purpose of separating $L_1$, $L_2$, and the deeper interior is that sensitivities only depend on the outer two layers, while the deeper points are only needed when restoring those layers after deletions.

\subsection*{Why the running time is near-linear}
The challenge is not recomputing one convex hull, but maintaining the entire peeling process efficiently. The proof strategy is amortized:
\begin{itemize}
    \item hull restoration is charged to points that move outward through convex layers,
    \item sensitivity updates are charged to newly activated points,
    \item and the total number of such activations is $O(n)$.
\end{itemize}
This yields an overall running time of $O(n \log n)$ with $O(n)$ space for the full peeling order.

\subsection*{Implementation takeaway}
For an implementation roadmap, the thesis logic naturally separates into:
\begin{itemize}
    \item a geometric kernel for predicates and polygon area,
    \item a verifiable brute-force oracle,
    \item a layered hull maintenance strategy,
    \item tangent-query support,
    \item and finally a dynamic deletion-only hull engine.
\end{itemize}
This progression mirrors the practical development path from a correctness-first oracle to the intended high-performance implementation.

